

Two EEG corpora acquired during sustained simulated-driving tasks were used, providing 104 complete recording sessions in total. The corpora differed in their acquisition montages and data provenance, allowing the proposed method to be evaluated under different hardware setups.

The public Cao2018 corpus contains 62 sessions from 27 participants who performed an approximately 90-minute lane-keeping task in an immersive driving simulator~\citep{cao2019multichannel}. The present study included the 58 sessions marked as complete and excluded the remaining four incomplete sessions. This study used the raw recordings. All 32 recorded electrophysiological channels (30 scalp electrodes and the A1/A2 mastoid reference channels) were supplied to the pipeline; the concurrently recorded vehicle-position channel was excluded. These 32 channels were Fp1, F3, Fp2, F4, FT7, FC3, C3, FC4, FT8, C4, TP7, CP3, TP8, CP4, T5, P3, O1, T6, P4, O2, F7, T3, Fz, FCz, Cz, Pz, Oz, F8, T4, CPz, A1, and A2.

The Internal corpus comprised 46 complete sessions recorded during a simulated lane-keeping task. EEG was acquired with a 128-channel EGI HydroCel system, with CPz as the reference. Although the original recordings used a high-density montage, the present analyses retained a predefined subset of 32 electrodes corresponding to conventional 10--20 scalp locations, following the subsampled 128-to-32-channel montage of~\citep{Hatlestad_2023_reliable}: Fp1, AF3, F3, Fp2, AF4, F4, FC5, FC1, C3, FC2, FC6, C4, CP5, CP1, CP2, CP6, P5, P3, PO3, O1, PO4, P4, P7, O2, F7, T7, Fz, Cz, Pz, Oz, F8, and T8.

We divided the 32 electrodes into four anatomical regions (frontal, central, parietal, and occipital) and further divided each region into left and right hemispheres, as per~\citep{black2012}, as shown in Figure~\ref{fig:eeg_map}. In the figure, the resulting eight regions (frontal left, frontal right, central left, central right, parietal left, parietal right, occipital left, and occipital right) are abbreviated FL, FR, CL, CR, PL, PR, OL, and OR, respectively. The figure shows the scalp locations of the 32 electrodes: the blue circles represent electrodes located in both the Internal and Cao2018 montages, the yellow circles represent electrodes located only in the Internal montage, and the red circles represent electrodes located only in the Cao2018 montage.

Of these eight regions, only the frontal, central, parietal, and occipital regions (left and right) were used in the analysis. The following electrodes were not assigned to any lateralized region and were likewise excluded from the analysis: F7, T7, F8, T8, Fz, FCz, Cz, CPz, Pz, and Oz.

\begin{figure}[ht]
    \centering
    \includegraphics[width=0.9\linewidth]{eeg_map}
    \caption{Scalp locations of the 32 electrodes used in this study, divided into eight anatomical regions (frontal left/right, central left/right, parietal left/right, and occipital left/right, abbreviated FL, FR, CL, CR, PL, PR, OL, and OR). Blue circles mark electrodes present in both the Internal and Cao2018 montages, yellow circles mark electrodes present only in the Internal montage, and red circles mark electrodes present only in the Cao2018 montage.}
    \label{fig:eeg_map}
\end{figure}
